%! Polar Function Graphs — Unit 7, Lesson 7.7 — Exit Ticket
\documentclass[10pt]{article}
\usepackage{precalculus-article}
\usepackage{precalculus-boxes}
\newcommand{\TallMath}[1]{$\displaystyle #1\rule[-1.4em]{0pt}{3.2em}$}

\begin{document}

\pageheader{Unit 7, Lesson 7.7}{Exit Ticket}
\namedateperiod

\vspace{0.12in}

\begin{notesbox}{Exit Ticket --- Answer independently}

\textbf{1.}\quad Complete the table for $r = 3\cos\theta$.

\vspace{4pt}
\renewcommand{\arraystretch}{1.7}
\begin{tabular}{|c|c|c|c|c|}
\hline
$\theta$ & $0$ & $\dfrac{\pi}{2}$ & $\pi$ & $\dfrac{3\pi}{2}$ \\
\hline
$r = 3\cos\theta$ & & & & \\[4pt]
\hline
\end{tabular}

\vspace{0.18in}
\textbf{2.}\quad Based on your table, at what value(s) of $\theta$ in $[0, 2\pi)$
does $r = 0$ for the function $r = 3\cos\theta$?

\vspace{6pt}
$\theta = $ \underline{\hspace{4cm}}

\vspace{6pt}
What does this mean for the polar graph at those angles?

\writeline

\vspace{0.16in}
\textbf{3.}\quad (Multiple Choice) Which statement \textbf{best} describes how $\theta$
and $r$ relate in a polar function $r = f(\theta)$?

\vspace{6pt}
\begin{enumerate}[label=(\Alph*), leftmargin=2em, itemsep=6pt]
  \item $\theta$ is the output (radius) and $r$ is the input (angle).
  \item $\theta$ is the input (angle) and $r$ is the output (radius).
  \item Both $\theta$ and $r$ are inputs to the function.
  \item Both $\theta$ and $r$ are outputs of the function.
\end{enumerate}

\vspace{4pt}
Answer: \underline{\hspace{1.5cm}}

\end{notesbox}

\end{document}
